\documentclass[whitelisted,showhints]{nhrproposal}

\newcommand\projecttitle{[Project title here]}
\newcommand\authors{[PI/PC of the proposal here]}
\newcommand\projectlogo{figures/some_logo.png} %comment out if no logo should be included

\begin{document}
\Titlepage

\ShortProjectDescription
\ShortSubprojectDescription{[Sub-project title]}
\ShortSubprojectDescription{[Sub-project title]}

\ComputeResources
\CodeOverview
%for one code
\begin{Codetable}{1}{0.6\linewidth}{0.3\linewidth}{0}{0} %first value: number of codes/programs, other values: column widths
%for two codes
%\begin{Codetable}{2}{0.5\linewidth}{0.2\linewidth}{0.2\linewidth}{0} %first value: number of codes/programs, other values: column widths
%for three codes
%\begin{Codetable}{3}{0.4\linewidth}{0.2\linewidth}{0.2\linewidth}{0.2\linewidth} %first value: number of codes/programs, other values: column widths
 Name           &  \\ \hline
 Version number &  \\ \hline
 Web page       & \url{...}   \\ \hline
 Citation or reference &   \\ \hline 
 Are you a developer, collaborator or end user of the program? &   \\ \hline
 Is the source code available? (publicly, upon request, no) &   \\ \hline
 Is it a commercially licensed software? If yes, do you already have a license that is useable at the requested NHR center?  &  \\ \hline
 Are there usage limitation for this software like the maximal number of simultaneous runs, license server,...?  &  \\ \hline
 How is the code parallelized (pure MPI, hybrid MPI/OpenMP, CUDA,...)? &  \\ \hline
 Which programming languages are used? (optional) &  \\ \hline
 Are there other further requirements (e.g. special libraries) for the program? (optional) &   \\ \hline
\end{Codetable}

\CodeBenchmarks
\Benchmark{Program A Name}
\BenchmarkDescription

\BenchmarkSystem
\begin{center}
\begin{tabular}{|G{5cm}|L{10cm}|}\hline
Cluster location &  \\ \hline
Cluster name & \\ \hline
CPUs per node & \\ \hline
CPU type & \\ \hline
Physical CPU-cores per CPU & \\ \hline
Main memory per node & \\ \hline
Interconnect & \\ \hline
Accelerators (such as GPUs) & \\ \hline
\end{tabular}
\end{center}

\BenchmarkData

\BenchmarkDiscussion

\Benchmark{Program B Name}
%\BenchmarkDescription
%
%\BenchmarkSystem
%\begin{center}
%\begin{tabular}{|G{5cm}|L{10cm}|}\hline
%Cluster location &  \\ \hline
%Cluster name & \\ \hline
%CPUs per node & \\ \hline
%CPU type & \\ \hline
%Physical CPU-cores per CPU & \\ \hline
%Main memory per node & \\ \hline
%Interconnect & \\ \hline
%Accelerators (such as GPUs) & \\ \hline
%\end{tabular}
%\end{center}

%\BenchmarkData

%\BenchmarkDiscussion

\ResourceJustification
%ResourceTable
%first argument: resource type, e.g, CPU, GPU, FPGA,...
%second argument: resource usage specification, i.e., subproject;run type;programs used;problem size;#runs;#steps per run; Wall time per step [hours];#CPU-cores per run
\ResourceTable{CPU}{
%subproject;run type;programs used;problem size;#runs;#steps per run; Wall time per step [hours];#CPU-cores per run
1;Preproc;program A;N=1000;100;20;3.4;128\\
1;Type1;program B;N=15000;1000;20;3.4;128\\
2;Type2;program C;N=100000;100;200;3.4;512\\
}

%alternative ResourceTable for molecular dynamics simulations
%\ResourceTableMD{CPU}{
%%subproject;run type;programs used;problem size;#runs;simulation time in mus; performance ns/day;#CPU-cores per run
%1;Type1;program B;N=100k;10;15;110.5;128\\
%}

\ResourceTable{GPU}{
%subproject;run type;programs used;problem size;#runs;#steps per run; Wall time per step [hours];#GPUs per run
1;Type1;program B;N=1000;100;20;3.4;4\\
}

%alternative ResourceTable for molecular dynamics simulations
%\ResourceTableMD{GPU}{
%%subproject;run type;programs used;problem size;#runs;simulation time in mus; performance ns/day;#CPU-cores per run
%1;Type1;program B;N=100k;10;15;110.5;4\\
%}


%manual resource table
%\begin{tabular}{|p{1.3cm}|p{2cm}|p{2cm}|p{2cm}|p{1.2cm}|p{1.2cm}|p{1.3cm}|p{1.2cm}|p{1.7cm}|} \hline
%    \rowcolor{Gray}    Sub-project & Run type & Programs used &  Problem size & \#runs & \#steps per run & Wall time per step [hours] & \#CPU-cores / run & Total \\ \hline
%  1 & Preproc & program A & $PS1$ & $R1$ & $S1$ & $W1$ & $C1$ & $R1\cdot S1\cdot W1\cdot C1\cdot 10^{-6}$ \\ \hline
%  ... & Type 1 & program B & $PS2$ & $R2$ & $S2$ & $W2$ & $C2$ & $R2\cdot S2\cdot W2\cdot C2\cdot 10^{-6}$ \\ \hline
%  TOTAL& - & - & -& -&-&-&-& Sum here \\ \hline 
%\end{tabular}


\ResourcePlan
%PlanTable
%first argument: table width in cm
%second argument: identifiers of resource types for sum columns, e.g. CPU, GPU, FPGA,...
%third argument: planed usage separated by sub projects, resource types and months
\PlanTable{14}{CPU,GPU}{ 
  Sub-project 1;
    CPU,0.1,1.2,3.3,1.4,1.5,1.6,1.7,1.8,1.9,2.10,2.11,2.12;
    GPU,1.1,1.2,3.3,1.4,1.5,1.6,1.7,1.8,1.9,2.10,2.11,2.12;
  \\
  Sub-project 2;
    CPU,1.1,1.2,3.3,1.4,1.5,1.6,1.7,1.8,1.9,2.10,2.11,2.12;
  \\
}


\SpecialRequirements

\References{references} %first argument is name of bibtex file

%\BenchmarkExample

\end{document}
